% !TEX TS-program = pdflatex
\documentclass[final,3p,times]{elsarticle}
\pdfminorversion=7

\usepackage{amsmath,amssymb}
\usepackage{booktabs}
\usepackage{tabularx}
\usepackage{listings}
\usepackage{url}
\usepackage[hidelinks]{hyperref}
\usepackage{xcolor}
\usepackage{placeins}

\lstset{
  basicstyle=\ttfamily\small,
  breaklines=true,
  breakatwhitespace=false,
  showstringspaces=false,
  columns=fullflexible
}

\bibliographystyle{zstar-elsarticle-num}
\journal{Computer Physics Communications}

\begin{document}

\begin{frontmatter}
\title{VCNEB: A calculator-agnostic Python toolkit for variable-cell nudged elastic band calculations}

\author{Xudong Zhu}
\address{VCNEB developers, China}

\begin{abstract}
Solid--solid phase transitions often involve coupled atomic rearrangements,
shear, and changes in the simulation cell.  A fixed-cell nudged elastic band
(NEB) calculation can therefore constrain the mechanism or assign an
ambiguous energetic cost to the transition.  We present VCNEB, a pure Python
toolkit that extends NEB to an explicit atomic--cell configuration space while
keeping the electronic-structure calculator behind an ASE-compatible contract.
The implementation uses fractional atomic coordinates and a deformation
gradient, converts calculator stress to a work-consistent cell force, and
reuses mature ASE optimizers for FIRE, BFGS, and LBFGS updates.  It provides
climbing images, pressure-dependent enthalpy paths, cell and atomic component
masks, mode-guided initialization, strict or projected mode constraints,
restartable trajectories, image-level parallel ownership, and structured
force/stress diagnostics.  VASP and ABACUS adapters are included.  Analytic
finite-difference tests, a coupled atomic--cell double-well, a fixed-cell ASE
comparison, and a release-and-refine mode workflow validate the implementation.
A mapped 12-atom HfO$_2$ tetragonal-to-polar-orthorhombic case demonstrates the
cluster workflow and exposes the diagnostics required before a production
barrier is reported.  VCNEB is intended as an open, reproducible software
foundation for variable-cell transition-state studies rather than as a claim
that the variable-cell NEB algorithm itself is new.
\end{abstract}

\begin{keyword}
variable-cell NEB \sep solid--solid phase transition \sep stress \sep
climbing image \sep ASE \sep VASP \sep ABACUS
\end{keyword}
\end{frontmatter}

{\bf PROGRAM SUMMARY}
\begin{small}
\noindent
{\em Program Title: VCNEB} \\
{\em CPC Library link to program files:} to be added under proof \\
{\em Developer's repository link:} \url{https://github.com/xdzhu/vcneb} \\
{\em Licensing provisions:} to be selected before release \\
{\em Programming language:} Python \\
{\em Nature of problem:} Calculation of minimum-energy paths and transition-state candidates for crystal transformations involving both atomic motion and changes of the simulation cell. \\
{\em Solution method:} VCNEB combines fractional atomic coordinates with a deformation-gradient cell representation, transforms calculator stress into a generalized cell force, applies NEB/CI-NEB projections in the resulting extended metric, and delegates numerical updates to established ASE optimizers. \\
\end{small}

\section{Introduction}
\label{sec:introduction}

The nudged elastic band method (NEB) is a standard approach for finding a
minimum-energy path between two local minima \cite{Henkelman2000}.  In a
crystal, however, the endpoints of a transformation may differ in lattice
vectors as well as atomic positions.  A fixed-cell path can then hide a shear,
volume change, or cell-mediated collective coordinate and can produce a
mechanism that is appropriate only for an imposed boundary condition.

Several variable-cell or solid-state NEB formulations already exist.  Qian
\emph{et al.} introduced a variable-cell NEB method (VC-NEB) for structural
phase transitions \cite{Qian2013}; generalized solid-state NEB (G-SSNEB) uses a
closely related extended configuration space \cite{Sheppard2012}; and finite
deformation NEB makes the work-conjugate stress measure explicit
\cite{Ghasemi2019}.  Implementations are also available in USPEX, TSASE,
VTST, FLOS, ABINIT, and commercial atomistic simulation environments.  These
efforts establish the physical and algorithmic basis of cell-aware NEB, but
their practical workflows differ in calculator support, distribution model,
restart semantics, and diagnostic detail.

The purpose of VCNEB is therefore not to claim the first variable-cell NEB
algorithm.  It is to provide a small, auditable, calculator-agnostic Python
layer that can be placed between modern electronic-structure engines and
reproducible path workflows.  The project uses ASE for structure objects and
optimizers, but keeps VASP- and ABACUS-specific input generation outside the
path mechanics.  The same design also separates three meanings of a specified
vibrational mode: bending an initial path, projecting optimization updates,
and restricting the path to a hard subspace.

The remainder of this paper follows five sections.  Section~\ref{sec:theory}
defines the extended coordinates, stress transformation, NEB forces, and mode
constraints.  Section~\ref{sec:software} describes the software architecture,
calculator contract, restart model, and command-line workflow.
Section~\ref{sec:examples} presents numerical verification, model results,
calculator smoke tests, and the HfO$_2$ case study.  Section~\ref{sec:conclusions}
concludes with availability, limitations, and the path toward a production
release.

\section{Theory}
\label{sec:theory}

\subsection{Extended configuration space}

Let $h$ be the $3\times3$ cell matrix whose rows are the Cartesian cell
vectors, and let $s_a$ be the fractional coordinate of atom $a$.  VCNEB uses
the row-vector convention
\begin{equation}
  h = h_0 F^{\mathsf T}, \qquad r_a=s_a h,
  \label{eq:cell_convention}
\end{equation}
where $h_0$ is the reference cell and $F$ is a deformation gradient.  An
image is represented by $(s,F)$, or by a length-like generalized coordinate
\begin{equation}
  x = \left(\operatorname{vec}(s h_0),
       \operatorname{vec}\left[\lambda_c(F-I)\right]\right).
  \label{eq:extended_coordinate}
\end{equation}
The scale $\lambda_c$ has units of length and makes atomic and cell
components numerically comparable.  The default is the cube root of the
reference-cell volume, while users can supply a physically motivated value.
The interpolation operates on unwrapped fractional coordinates and linearly
interpolates $F$.  A positive determinant is required for every endpoint and
interpolated cell; global cell rotation is removed during the default endpoint
alignment.

\subsection{Enthalpy gradient and cell force}

Under hydrostatic pressure $P$, the path objective is the enthalpy
\begin{equation}
  H(s,F)=E(s,F)+P\Omega(F),
  \qquad \Omega=\det(h).
  \label{eq:enthalpy}
\end{equation}
Forces returned by a calculator are Cartesian forces $f_a=-\partial E/\partial
r_a$, and the fractional-coordinate force used by VCNEB is
\begin{equation}
  f_s = f_{\rm cart} h^{\mathsf T}.
  \label{eq:fractional_force}
\end{equation}
With the ASE stress convention and positive $P$ denoting compression, define
\begin{equation}
  W=-\Omega(\sigma+PI), \qquad
  G=\operatorname{solve}(F,W^{\mathsf T})^{\mathsf T}.
  \label{eq:cell_force}
\end{equation}
Here $G$ is conjugate to the components of $F$.  VCNEB maps $f_s$ and $G$ to
the generalized coordinate in Eq.~\eqref{eq:extended_coordinate}; the cell
block is divided by $\lambda_c$.  The signs and transposes in
Eq.~\eqref{eq:cell_force} are tested by central finite differences for
nonorthogonal cells and nonzero pressure.  The code rejects a variable-cell
run when an attached calculator does not expose stress.

\subsection{NEB and climbing-image forces}

For internal image $i$, let $\hat\tau_i$ be an energy-weighted tangent in the
extended coordinate space, and let $g_i$ be the true generalized enthalpy
force after component or mode projection.  With neighboring distances
$d_i^-$ and $d_i^+$, the non-climbing NEB force is
\begin{equation}
  F_i^{\rm NEB} = g_i- (g_i\cdot\hat\tau_i)\hat\tau_i
       + k_i(d_i^+-d_i^-)\hat\tau_i .
  \label{eq:neb_force}
\end{equation}
For the highest-energy interior image, the climbing-image force is
\begin{equation}
  F_i^{\rm CI}=g_i-2(g_i\cdot\hat\tau_i)\hat\tau_i,
  \label{eq:ci_force}
\end{equation}
so that the spring force is removed and the tangential physical force is
reversed.  Endpoints are never updated.  The primary optimizer criterion is
the maximum norm of the three-component blocks of the generalized force; the
software also reports maximum raw atomic force, stress, cell force, path
spacing, and a true/spring/NEB decomposition for every image.

\subsection{Modes and directional constraints}

Let $v_m$ be a normalized mode in the same extended metric as $x$.  A
mode-guided initial path is generated as
\begin{equation}
  x_i=x_i^{\rm lin}+A\,b(i/(N-1))v_m,
  \qquad b(0)=b(1)=0,
  \label{eq:guided_path}
\end{equation}
and is subsequently optimized without a constraint.  This is distinct from
projected dynamics, which applies the orthogonal projector
$\Pi=VV^{\mathsf T}$ to the true force and optimizer update, and from a hard
subspace path, for which the endpoints and all images are restricted to an
affine subspace.  Multiple nonorthogonal modes are orthogonalized by SVD.
Atomic masks and cell masks are applied before the mode projector, and a
diagnostic reports partially clipped or fully inactive basis columns.

The resulting constrained saddle is a saddle on the selected subspace, not
necessarily a first-order saddle of the full potential-energy surface.  The
recommended workflow is consequently to use a mode or directional constraint
to search for a mechanism, then release the constraint and refine the same
chain in the full extended space.

\section{Software}
\label{sec:software}

\subsection{Architecture and public interfaces}

The core package depends only on NumPy and ASE.  The central object,
\texttt{VCNEB}, exposes an ASE-compatible optimizable object whose coordinates
contain all internal images.  ASE's FIRE, BFGS, or LBFGS classes perform the
actual numerical update; VCNEB supplies the generalized forces and handles
cell validity, constraints, climbing images, and endpoint ownership.

The calculator contract is deliberately small: each image must provide
\texttt{get\_potential\_energy()}, \texttt{get\_forces()}, and
\texttt{get\_stress(voigt=False)}.  A preflight validator checks declared
capabilities and, where requested, unique per-image directories.  Runtime
errors include the image index, calculator type, working directory, launcher,
and original exception.  The VASP adapter copies source input settings while
forcing static image evaluations; the ABACUS adapter enforces
\texttt{cal\_force=1}, \texttt{cal\_stress=1}, and \texttt{out\_stru=1}.

\subsection{Interpolation, restart, and provenance}

Structures must have equal atom counts and elemental composition.  The default
identity order remains backward compatible, while an explicit permutation or
an element-grouped geometry mapping can be requested and saved as a report.
The interpolator supports periodic minimum-image fractional displacements,
cell alignment, linear or logarithmic-strain cell paths, user callbacks,
positive-volume guards, and exact endpoint preservation.  Every
optimizer callback can write a complete chain snapshot and a flat ASE
trajectory.  Resume logic counts only complete groups of $N$ image frames, so
a truncated tail does not corrupt the restart chain.  Example drivers record
the source commit, node check, calculator settings, work directory, summary
JSON, and human-readable log.

\subsection{Typical use}

The following fragment is sufficient for a custom ASE calculator:
\begin{lstlisting}[language=Python]
from vcneb import interpolate_vcneb, run_vcneb

images = interpolate_vcneb(initial, final, n_images=9, mic=True)
for image, calculator in zip(images, calculators):
    image.calc = calculator

chain, optimizer = run_vcneb(
    images, k=0.10, climb=True, optimizer="FIRE",
    fmax=0.05, steps=300,
    trajectory="vcneb.traj", snapshot_dir="snapshots",
)
print(chain.barrier())
print(chain.path_diagnostics())
\end{lstlisting}

The command-line ABACUS driver accepts endpoint structures, pseudopotential
and orbital mappings, an MPI launcher, and the optimizer settings.  The VASP
driver follows the same path mechanics while copying INCAR, KPOINTS, and
POTCAR from an endpoint template.  These adapters are optional runtime
configurations rather than core dependencies.

\section{Examples}
\label{sec:examples}

All numerical examples reported here are tied to a repository commit and an
input/output manifest.  Long electronic-structure jobs were executed on the
shared 40-core nodes through gateway 235; no DFT or NEB production calculation
was run on the workstation.

\subsection{Differential tests and fixed-cell reduction}

Central finite-difference tests on a nonorthogonal cell give maximum absolute
errors of $3.18\times10^{-11}$ for the fractional force and
$5.66\times10^{-11}$ for the cell force in the current regression suite.  A
separate six-direction physical strain scan gives $3.16\times10^{-11}$ eV at
zero pressure and $1.35\times10^{-10}$ eV at $0.5$ GPa, with the stable
central-difference step $10^{-6}$.  The pressure-dependent enthalpy derivative
and the cell/atomic mask behavior are also covered.  In a coupled atomic--cell double-well with a known saddle,
CI-VCNEB recovers the $0.25$ eV barrier, the saddle at
$(q_x,\epsilon_{xx})=(0.5,0.125)$, a generalized residual below $0.01$ eV/Å,
and negative local tangent curvature.

Setting the cell mask to zero reduces VCNEB to a fixed-cell NEB in the same
extended coordinate convention.  Using an identical seven-image analytic
path, VCNEB gives a barrier of $0.28116144$ eV, compared with $0.28117077$ eV
from ASE CINEB; the absolute difference is $9.33\times10^{-6}$ eV.  The
remaining difference is attributable to optimizer stopping histories rather
than a different physical calculator.

\subsection{Mode-guided search and release-and-refine}

The analytic workflow starts with a transverse mode perturbation, performs a
projected collective-mode search, and then releases the constraint.  On the
cluster, the constrained stage converged in 12 FIRE steps with a barrier of
0.286000 eV and a maximum generalized force of 0.009891 eV/Å.  The released
full-space refinement converged in 37 steps with a barrier of 0.250044 eV and
a maximum generalized force of 0.008655 eV/Å.  The 0.035956 eV decrease is a
direct demonstration that a constrained mechanism barrier must not be
reported as the unconstrained transition barrier.

\subsection{Calculator smoke tests and HfO$_2$ diagnostic case}

The VASP adapter was tested on a small GaN cell and returned finite energy,
forces, and stress using VASP 6.3.2.  The ABACUS adapter was tested on the
mapped HfO$_2$ endpoint using Dojo-NC-FR pseudopotentials, 60 Ry, 1$\times$1$\times$1
k points, and an SCF threshold of $10^{-6}$.  The single-image result was
finite and passed the stress validator.  These are calculator smoke tests,
not convergence claims for a phase-transition barrier.

The material fixture contains 12 atoms with an explicit T-to-PO atom mapping.
Independent loose endpoint relaxations reached maximum atomic forces below
0.1 eV/Å in the recorded trial.  A seven-image variable-cell path was then
run as an engineering trial.  A 40-step FIRE branch and an LBFGS recovery
branch both failed to converge; a conservative FIRE retry with
\texttt{maxstep=0.05} still increased the residual from 0.617654 to
0.725727 eV/Å in seven complete steps.  A zero-step static evaluation of its
latest chain found image 3 as the highest enthalpy image (relative enthalpy
0.305708 eV), while image 2 carried the largest raw interior force
(0.674825 eV/Å) and stress component (0.016611 eV/Å$^3$).  The mismatch is
useful evidence that the current linear/restarted path requires a
mechanism-aware improvement; none of these HfO$_2$ values is reported as a
production transition barrier.

\subsection{BaTiO$_3$ endpoint and variable-cell diagnostic}

As a compact perovskite test, we also prepared a five-atom cubic-to-tetragonal
BaTiO$_3$ pair.  The endpoint cells differ in the tetragonal axis and the
oxygen order is nontrivial: the identity path reaches a minimum periodic
interatomic distance of only $0.072$~\AA, whereas the element-grouped mapping
\texttt{[0,1,3,2,4]} gives $1.347$~\AA.  This is a small but concrete example
of why endpoint mapping must be audited before a calculator is attached.

Independent variable-cell endpoint trials with the shared PBE/LCAO ABACUS
setup produced volumes of $90.8105$ and $89.2083$~\AA$^3$ for the cubic and
tetragonal endpoints, respectively.  A subsequent seven-image logarithmic
strain VCNEB trial used FIRE with \texttt{maxstep=0.02} for 40 steps.  Its
maximum generalized residual was $3.380$~eV/\AA; all interior images were below
the cubic endpoint, giving a formal forward barrier of zero and a reaction
enthalpy of $-13.582$~eV.  This is an engineering and adapter diagnostic, not
a physical BaTiO$_3$ barrier: the tetragonal endpoint itself remained just
above the force target, the interior residual was large, and the low-cost
electronic setup changes the equilibrium volume substantially.  The complete
summary, endpoint data, logs, and initial-path reports are included under
\texttt{outputs/batio3/}.

\begin{table}[!htbp]
\centering
\caption{Representative verification and diagnostic results.}
\label{tab:examples}
\begin{tabularx}{\linewidth}{l>{\raggedleft\arraybackslash}X>{\raggedleft\arraybackslash}X}
\toprule
Case & Main result & Status \\
\midrule
Finite difference & $3.16\times10^{-11}$ / $1.35\times10^{-10}$ & Passed \\
Analytic CI double well & $0.250000$ eV barrier & Passed \\
ASE fixed-cell comparison & $9.33\times10^{-6}$ eV difference & Passed \\
Release and refine & $0.286000\rightarrow0.250044$ eV & Passed \\
HfO$_2$ variable-cell trial & residual $0.622$ eV/Å & Diagnostic only \\
BaTiO$_3$ relaxed-endpoint trial & residual $3.380$ eV/Å & Diagnostic only \\
\bottomrule
\end{tabularx}
\end{table}

\FloatBarrier
\section{Conclusions/Availability}
\label{sec:conclusions}

VCNEB provides a compact pure Python implementation of a variable-cell NEB
workflow with explicit cell coordinates, stress-derived generalized forces,
climbing images, mode and direction controls, calculator adapters, restart
snapshots, and per-image diagnostics.  The numerical tests establish the
force transformation and fixed-cell reduction, while the release-and-refine
example establishes the intended semantics of constrained mechanism searches.
The HfO$_2$ and BaTiO$_3$ trials demonstrate that real material calculations
can be launched, resumed, audited, and stopped without confusing an
unconverged engineering path with a physical barrier.  The BaTiO$_3$ case also
shows why endpoint mapping and independent stress-aware relaxation belong in
the standard workflow.

The repository is available at
\url{https://github.com/xdzhu/vcneb}.  It contains the source package, tests,
examples, endpoint structures and mapping files, calculation manifests,
cluster logs, and this manuscript draft.  The current development version is
not yet a production release.  Before publication, we will add a clean
installation test, continuous integration for the non-DFT suite, a second
material case, stricter endpoint and electronic-structure convergence, and a
work-conjugate finite-deformation stress option.  A final release should also
freeze the license and citation metadata, provide an API tutorial, and include
the raw calculator inputs needed to reproduce every paper result.

The principal limitations are the requirement for a one-to-one endpoint atom
mapping, the current linear/logarithmic deformation-gradient interpolation,
the current hydrostatic-pressure formulation, and the absence of a converged
production HfO$_2$ or perovskite barrier.  Large reconstructive transitions may require
multiple initial paths, finite-wavevector supercells, magnetic or electronic
state controls, and a stress measure matched to the chosen finite-deformation
coordinate.  These are scientific controls on the calculation, not details
that can be hidden behind a generic optimizer interface.

\bibliography{vcneb}
\end{document}
